% \iffalse meta-comment
%
% hit-thesis-paragraph.dtx 是段落与行距
%
% \fi
%
% \section{段落与行距}
%
% ^^A ======================================================================
% ^^A Module thesis-paragraph: 段落首行缩进与标点处理（hit-thesis）
% ^^A ======================================================================
%    \begin{macrocode}
%<*thesis-paragraph>
\ctexset{%
  punct=quanjiao,
  space=auto,
  autoindent=true}
%    \end{macrocode}
%
% 调整默认列表环境间的距离，以符合中文习惯。机制在
% \file{src/utils/hit-list.dtx}，这里只给取值。
%    \begin{macrocode}
\bool_gset_true:N \g__hit_list_nosep_bool
\clist_gset:Nn \g__hit_list_itemize_clist
  { leftmargin=0em,itemsep=0em,partopsep=0em,parsep=0em,topsep=0em,itemindent=3em }
\clist_gset:Nn \g__hit_list_enumerate_clist
  { leftmargin=0em,itemsep=0em,partopsep=0em,parsep=0em,topsep=0em,itemindent=3.5em }
\hit@setup@list@spacing
%    \end{macrocode}
%
% \changes{v3.2a}{2026/08/07}{列表间距的三行迁进 \file{src/utils/hit-list.dtx}，
%   与 \env{publist} 同处；这里只留调用与正文首行缩进}
% \changes{v3.2a}{2026/08/15}{首行缩进由 \texttt{2em} 改成 \texttt{2\cs{ccwd}}，
%   量的从西文 em 换成汉字宽}
% 正文首行缩进。规范要的是「空两个汉字」，所以量 \cs{ccwd}（\pkg{ctex} 的汉字宽）
% 而不是 \texttt{em}（当前\emph{西文}字体的 em）。两个数在这里恰好相等，
% 宋体一个字与 Times 的一个 em 都等于字号，所以一直没出岔子。
%
% 但拴在西文字体上会出事：西文字体加载失败时退回 \cs{nullfont}，它的
% \cs{fontdimen}~6 是零，\texttt{2em} 静默变成 0\,pt，正文首行缩进整个消失，
% 日志里只有一句 \texttt{not loadable} 混在一堆 \pkg{fontspec} 报错里。
% \pkg{newtx} 没装（\option{fontset} 取 \texttt{fandol}、\texttt{ubuntu} 时
% 西文走 \texttt{TeXGyreTermesX}，那个字体是 \pkg{newtx} 带的）就会这样。
% report 类那边 \issue[201] 已经改成 \texttt{2\cs{ccwd}} 了，这里跟上。
%
% 新版面不走这一句：\cs{hit@format@apply@body:} 会用字符网格的步进重设
% \cs{parindent}，那才是 Word 的算法（范例存的是
% \texttt{firstLineChars="200"} 加 \texttt{firstLine="498"} 缇，
% 498/20/2 = 12.45\,bp 一个字符，等于字号加字距，不是字号）。
%    \begin{macrocode}
\dim_set:Nn \parindent {2\ccwd}
%</thesis-paragraph>
%    \end{macrocode}
%
% \Finale
